\documentclass[conference]{IEEEtran}

\usepackage{amsmath}
\usepackage{graphicx}
\usepackage{hyperref}
\hypersetup{colorlinks=true, urlcolor=blue}

% Make section titles visually bolder
\makeatletter
\def\section{\@startsection{section}{1}{0mm}
  {-2ex plus -.5ex minus -.2ex}
  {1ex plus .2ex}
  {\normalfont\large\bfseries}}
\makeatother

\title{
Maternal Mental Health Risk Assessment:\\
A Fuzzy Logic-Based Machine Learning Approach for
Postpartum Depression Screening in Nepal
}

\author{
    \IEEEauthorblockN{Rajesh Kumar Dhimal}
    \IEEEauthorblockA{
        Student ID: 250072, CUID: 16542745\\
        MSc. Data Science and Computational Intelligence\\
        Softwarica College of IT and E-commerce, Coventry University\\
        250072@softwarica.edu.np
    }
}

\begin{document}
\maketitle
\vspace{-14pt}

\begin{abstract}
Postpartum depression (PPD) remains under-recognized in Nepal due to stigma, limited screening tools, and the lack of microdata. This proposal presents a compact framework integrating fuzzy logic and machine learning, supported by a rigorously constructed synthetic population. The goal is to provide an interpretable, context-aware screening method for early PPD risk detection in resource-limited settings.
\end{abstract}

\vspace{6pt}  % adjust value (6pt, 8pt, 10pt...)
\begin{IEEEkeywords}
Fuzzy Logic, Maternal Mental Health, Postpartum Depression, Machine Learning, Synthetic Data, Nepal
\end{IEEEkeywords}

\vspace{-6pt}
\section{Introduction}
Maternal mental health issues, especially PPD, pose significant public health challenges in Nepal. Traditional scoring methods struggle to represent the uncertainty and overlapping risk patterns inherent in mental health states. A fuzzy logic-based model, refined using machine learning, offers a flexible and explainable alternative suited to Nepal’s socio-cultural diversity.

\vspace{-4pt}
\section{Aims and Objectives}
\begin{itemize}
    \item Develop a fuzzy logic screening model for postpartum depression risk.
    \item Use machine learning to benchmark, refine, and support fuzzy inference.
    \item Create a validated synthetic maternal health dataset grounded in national evidence.
\end{itemize}

\vspace{-6pt}
\section{Data Synthesis}
A synthetic dataset was generated after an extensive review of statistical metrics, distributions, and relationships observed in:
\begin{itemize}
    \item Nepal Demographic and Health Survey (NDHS) 2022,
    \item Nepal Population Census 2021,
    \item WHO, “Maternal Mental Health: Evidence and Practice,” 2021.
\end{itemize}

A complete simulated population representing Nepal’s mother-age 
demographic was first created using conditional, evidence-based rules. 
This population data was passed through the synthetic data quality tests 
(e.g, distributional similarity, cross-tab consistency, correlation 
preservation). A final working sample of 20,000 records was drawn 
using simple random sampling to ensure 
representativeness and reliability.

\vspace{-6pt}
\section{Methodology}
\begin{itemize}
    \item \textbf{Fuzzy Feature Engineering:} Construction of fuzzy-enhanced variables (support level, stress intensity, sleep quality, household strain, etc.) using membership functions that capture gradual transitions rather than strict thresholds.
    \item \textbf{Fuzzy Logic Engine:} Expert-informed fuzzy rules infer PPD risk levels based on socio-demographic, psychosocial, and maternal health indicators.
    \item \textbf{Machine Learning Model Selection:}
    \begin{itemize}
        \item \textbf{Primary Algorithm:} Fuzzy K-Nearest Neighbors (Fuzzy KNN) for smooth decision boundaries and robustness to overlapping risk patterns.
        \item \textbf{Secondary Algorithms:}  
              (1) Fuzzy Decision Tree,  
              (2) Fuzzy C-Means Clustering + Classification,  
              (3) Fuzzy Naive Bayes.
    \end{itemize}
    The best-performing model will be selected based on accuracy, interpretability, and consistency with expected PPD prevalence patterns.
    \item \textbf{Validation:} Predicted risk distributions compared with published PPD prevalence ranges and evaluated across provinces and demographic strata.
\end{itemize}

\vspace{-6pt}
\section{Expected Outcomes}
\begin{itemize}
    \item A Nepal-specific fuzzy logic PPD screening tool.
    \item Machine learning insights into key maternal mental health predictors.
    \item A validated, documented synthetic dataset for future research.
\end{itemize}

\vspace{-6pt}
\section{Ethical Considerations}
No real individuals’ data are used; all records are fully synthetic and derived from aggregate public sources, ensuring zero privacy risk.

\vspace{-4pt}
\section{Conclusion}
Integrating fuzzy logic, machine learning, and synthetic epidemiological modeling provides a practical, interpretable, and scalable pathway for early maternal mental health risk screening in Nepal.

%\bibliographystyle{IEEEtran}
%\%begin{thebibliography}{1}
%\bibitem{NDHS2022}
%Nepal Demographic and Health Survey 2022.
%\bibitem{Census2021}
%Nepal Population Census 2021.
%\bibitem{PPDMeta}
%Bista B., et al., ``Postpartum depression in Nepal,'' Asian J. Psychiatry, 2022.
%\bibitem{WHO2021}
%WHO, ``Maternal Mental Health: Evidence and Practice,'' 2021.
%\end{thebibliography}

\end{document}
